\chapter{知识复用的判断与选择}
\label{chap:ket-reuse-selection}

花卉识别系统准备进入新地区，目标端只能标注少量照片。我们可以取得一个物种分类模型、一个病害识别模型和一个花朵定位模型，但还不知道哪段经验能帮助当前目标。来源任务的准确率衡量它原来回答问题的能力，照片风格的相似度描述输入关系；两者都没有直接回答它能否减少目标错误。即使某个来源让训练前期很快达到可用准确率，继续训练后也可能不如随机初始化。另一些来源单独使用收益不大，组合后却可能补齐彼此遗漏的输入区域。

要作出选择，我们需要把这几个疑问分开。收益判断先规定相对于什么对照、改善哪个量才算有帮助；适用性判断检查来源保留的信息能否满足目标要求；来源选择再依据这些证据和实际预算决定使用哪些内容。第~\ref{chap:ket-overview}~章已区分知识载体与任务能力：保存了参数、表示或样本，并不等于在给定接口下能够完成目标任务。本章因此判断具体的复用对象及其使用方式，不给整个来源模型一个永久有效的通用性分数。

沿用前章的记号，$\mathcal P_S$与$\mathcal P_T$分别描述来源和目标的输入、标签联合分布，$\ell_T$为目标预测错误计分。目标标签可以用于拟合模型，也可以用于比较来源，后文将分别记录这两种用途。章内的小规模数值、假设成绩和花卉任务关系均为\term{教学构造}（teaching construction）：人为规定输入与结果，以便逐步复算判断，不作为论文实验或真实系统测量。文献中的方法和研究结论另注明出处。

\section{知识复用的收益判断}
\label{sec:ket-reuse-benefits}

在新地区的同一批目标照片上，复用方案和从头训练可能给出不同预测。解释这个差异以前，要检查二者是否得到相同的目标证据和可比较的训练机会。这里称不使用待检验来源知识、在约定信息条件下完成同一目标任务的方案为\term{无迁移对照}（no-transfer baseline）。对照应有合理的训练与调参程序；若复用方案多用了标签、更大的网络或更多搜索，观察到的差异就混入了这些因素的作用。

比较条件随所问的问题确定。要隔离来源初始化的作用，应尽量保持目标网络、训练数据和优化搜索规则一致，只改变参数起点。要在设备预算内选择完整方案，则可以比较冻结表示、全量微调和从头训练各自在该预算下经合理搜索得到的配置，并报告它们不同的更新范围与资源占用。后一种比较回答的是方案收益，不把所有差异都归因于初始化。

\subsection{预测效果收益的判断}
\label{sec:ket-reuse-prediction}

先考察训练终点的预测效果。令学习程序$\mathcal A$使用来源知识$K$和目标训练集$S_T$得到$f_K=\mathcal A(K,S_T)$；在其余规则相同、来源通道关闭时，得到$f_0=\mathcal A(\varnothing,S_T)$。前章已用风险描述未来目标输入上的平均损失。我们把无迁移风险减去复用风险，使正的差值对应错误减少：
\begin{equation}
  \begin{aligned}
    R_T(f)&=\mathbb E_{(\bm{x},y)\sim\mathcal P_T}
      [\ell_T(f(\bm{x}),y)],\\
    \Delta_T(K)&=R_T(f_0)-R_T(f_K).
  \end{aligned}
  \label{eq:ket-reuse-risk-gain}
\end{equation}
式中第一行的期望对目标输入与标签取平均，第二行比较两个已得到的预测器。在损失越小越好的约定下，$\Delta_T(K)>0$称为\term{正迁移}（positive transfer），表示使用来源后目标风险下降；$\Delta_T(K)<0$称为\term{负迁移}（negative transfer），表示使用来源后风险上升。若要评价学习程序而非这一次训练结果，还需对训练集抽样、初始化等随机因素取期望，并用重复运行估计。一次正差值不能单独证明平均收益为正。

Wang等以同一学习算法有无来源信息时的目标风险差定义负迁移，并分析算法、来源与目标的关系及目标标签量对结果的影响。式~\eqref{eq:ket-reuse-risk-gain}~采用风险下降量，符号与该文采用的负迁移间隙相反。如果复用算法输给另一种更强的无迁移算法，这仍是实际选型的重要证据；但算法也变了，尚不能据此把全部差距解释为来源知识的作用\scite{ket-wang2019-negative-transfer}

总体风险未知时，我们只能用有限目标样本估计。对不参与拟合、按目标抽样协议取得且非空的样本集$V_T$，将每个样本的目标损失相加，再除以样本数，得到
\[
  \hat R_{V_T}(f)
  =\frac{1}{|V_T|}
    \sum_{(\bm{x},y)\in V_T}\ell_T(f(\bm{x}),y).
\]
零一损失在预测错误时计1、正确时计0，其平均就是错误率，因此准确率增量等于错误率下降量。如果病害漏检更昂贵，应使用体现这种成本的损失，或另报病害阳性样本中被识别出的比例，即召回率。整体准确率不能替代这个要求，少数样本造成的差异也要连同估计不确定性解释。

表~\ref{tab:ket-reuse-checkpoints}~把比较固定为两物种识别：三种方案使用相同的目标网络结构、训练照片和100张验证照片，超参数搜索次数也相同。冻结方案保留来源表示，只训练新预测头；全量微调从来源参数起步更新整个网络；从头训练从随机参数起步。表中列出一次假设运行在第10步与第100步的错误数，检查点事先确定。读表时比较同一列；步数相同并不代表计算量相同，资源收益还需单独核算。

\begin{table}[htbp]
  \centering
  \small
  \begin{tabularx}{\linewidth}{Xrr}
    \toprule
    方案 & 第10步错误数 & 第100步错误数 \\
    \midrule
    冻结来源表示，训练新头 & 18 & 16 \\
    来源初始化，全量微调 & 12 & 14 \\
    随机初始化，从头训练 & 30 & 10 \\
    \bottomrule
  \end{tabularx}
  \caption{两物种识别的教学检查点。每个错误数的分母都是同一组100张验证照片；比较同一列得到该训练进度上的错误率差，不能由前一列外推后一列。这些数值不是实测学习曲线。}
  \label{tab:ket-reuse-checkpoints}
\end{table}

第10步全量微调相对于从头训练少错18张，错误率下降$(30-12)/100=0.18$，即18个百分点。到第100步，差值变为$(10-14)/100=-0.04$，即负4个百分点。若只能运行10步，早期优势可能满足上线要求；若预算允许运行100步，按早期结果锁定方案就会选错。这个构造只说明训练进度可以改变收益方向。微调后段错误增加的原因尚未给出，不能由两个检查点断定发生了过拟合或学习率失配。

冻结与微调的差别还决定结果能解释什么。冻结来源表示后训练新头，其成绩表明已有信息能否被当前预测头读出；微调成绩还包含改变表示的作用；从头训练给出没有该来源的学习结果。要判断来源训练是否改善了固定特征，可以补充冻结随机表示的对照；要判断初始化收益，则比较相同结构下的来源初始化与随机初始化。第~\ref{chap:ket-evaluation}~章说明如何配对重复运行、构造不确定性区间并预先规定训练终点。

\subsection{样本效率收益的判断}
\label{sec:ket-reuse-sample-efficiency}

如果上线要求已经规定准确率至少90\%，新的问题就是达到该效果要标注多少照片。对标签预算$n$，记按同一采样协议训练的预测器为$f_{K,n}$。把不同样本预算与所得效果对应起来，形成\term{学习曲线}（learning curve），这里用它观察增加目标标签时效果怎样改变。设可接受风险为$r_\star$，$N$是事先规定、实际测量的有限标签预算集合。只在已测点中寻找风险不高于阈值的最小预算：
\begin{equation}
  n_\star(K;r_\star)
  =
  \begin{cases}
    \min\{n\in N:\hat R_{V_T}(f_{K,n})\leq r_\star\},
      & \text{若集合非空},\\
    +\infty, & \text{若集合为空}.
  \end{cases}
  \label{eq:ket-reuse-label-threshold}
\end{equation}
取最小值是因为当前要固定效果、比较标签消耗，不是在固定标签下比较最高准确率。集合为空时的$+\infty$仅表示已测预算内没有过线，应同时报告最大已测预算，不能据此断言任何标签量都无效。式~\eqref{eq:ket-reuse-label-threshold}~得到的是经验到达量，依赖划分、类别构成、训练程序与阈值，不能单独证明总体风险达标。若重复运行存在波动，还须预先规定按平均风险、保守区间还是连续若干预算点判断过线。

图~\ref{fig:ket-reuse-learning-curves}~规定四个训练预算20、40、80、160。复用方案在同一100张验证照片上分别答对82、88、91、92张，从头训练分别答对65、76、85、91张；各预算都覆盖全部类别，较小训练集包含在较大训练集中。沿90\%水平线读图，复用与从头训练的首个已测达标预算分别为80与160，少用80个训练标签。沿80标签的竖线比较，二者准确率则相差6个百分点。这两次读取分别固定效果和固定标签，回答不同问题。

\begin{figure}[htbp]
  \centering
  % Teaching construction; no measured results.
% Training labels: 20, 40, 80, 160.
% Correct predictions out of 100: reuse 82, 88, 91, 92;
% scratch 65, 76, 85, 91. No interpolation claims or uncertainty estimates.
% Coordinates: x = labels * 0.045 cm; y = (accuracy - 60) * 0.12 cm.
\begin{tikzpicture}[
  x=1cm,y=1cm,font=\figurefont\footnotesize,
  line width=0.8pt,
  axis/.style={draw=BookInk,-{Latex[length=1.8mm]}},
  grid/.style={draw=BookRule,line width=0.5pt},
  reuse/.style={draw=BookTeal,line width=1.1pt},
  scratch/.style={draw=BookBlue,dashed,line width=1.1pt}
]
  \foreach \v/\yy in {60/0,70/1.2,80/2.4,90/3.6,100/4.8} {
    \draw[grid] (0,\yy) -- (7.55,\yy);
    \node[anchor=east,text=BookInk] at (-0.12,\yy) {\v};
  }
  \draw[BookAmber,densely dotted,line width=0.9pt]
    (0,3.6) -- (7.55,3.6);
  \draw[axis] (0,0) -- (7.85,0);
  \draw[axis] (0,0) -- (0,5.05);
  \foreach \n/\xx in {20/0.9,40/1.8,80/3.6,160/7.2} {
    \draw[BookInk] (\xx,0) -- (\xx,-0.08);
    \node[anchor=north,text=BookInk] at (\xx,-0.12) {\n};
  }
  \node[text=BookInk] at (3.9,-0.7) {目标训练标签数（个）};
  \node[rotate=90,text=BookInk] at (-0.85,2.4) {准确率（\%）};
  \draw[reuse] plot coordinates
    {(0.9,2.64) (1.8,3.36) (3.6,3.72) (7.2,3.84)};
  \foreach \xx/\yy in {0.9/2.64,1.8/3.36,3.6/3.72,7.2/3.84} {
    \fill[BookTeal] (\xx,\yy) circle[radius=2.2pt];
  }
  \draw[scratch] plot coordinates
    {(0.9,0.6) (1.8,1.92) (3.6,3.0) (7.2,3.72)};
  \foreach \xx/\yy in {0.9/0.6,1.8/1.92,3.6/3.0,7.2/3.72} {
    \draw[BookBlue,fill=white,line width=0.8pt]
      (\xx-0.06,\yy-0.06) rectangle (\xx+0.06,\yy+0.06);
  }
  \draw[reuse] (0.6,5.4) -- (1.25,5.4);
  \fill[BookTeal] (0.925,5.4) circle[radius=2.2pt];
  \node[anchor=west,text=BookInk] at (1.35,5.4) {复用};
  \draw[scratch] (3.0,5.4) -- (3.65,5.4);
  \draw[BookBlue,fill=white] (3.265,5.34) rectangle (3.385,5.46);
  \node[anchor=west,text=BookInk] at (3.75,5.4) {从头训练};
  \node[anchor=west,text=BookAmber] at (5.5,5.4) {阈值90\%};
\end{tikzpicture}

  \caption{标签预算与准确率的教学学习曲线。圆点表示复用，空心方点表示从头训练；90\%水平参考线用于读取首个已测达标预算80与160。连线不提供中间预算的测量结果。横轴只计训练标签，验证标签仍须计入开发总预算。}
  \label{fig:ket-reuse-learning-curves}
\end{figure}

假如到160标签仍没有过线，应报告“在已测预算内未达到，最大已测预算为160”，或按式~\eqref{eq:ket-reuse-label-threshold}~记为$+\infty$，不能把160写成到达量。未测预算可能因随机波动达标，因此也不能据此断定实际所需标签必然超过160。线段与阈值的交点不是实测的精确标签需求。比较病害任务时还要列出每类数量和标注粒度：完全没有阳性样本的80条记录与覆盖各病害类别的80条记录，提供的监督并不相同；一张照片的类别标签与逐像素掩码也不具有相同制作成本。

\begin{samepage}
训练标签的差值还没有包括来源选择。假设复用方案用了80个适应标签、30个与之不重叠的筛选标签及100个验证标签，开发合计210个；从头训练用了160个训练标签和同样100个验证标签，合计260个。完整开发阶段少用的是50个标签，而非80个。如果筛选的30个标签后来也进入适应，应按唯一标注对象计费，同时登记它的两种用途。已经用于拟合的数据不能因为先前被称为筛选集就继续承担独立验证；最终测试标签另外记账。

\end{samepage}

有些方案还读取大量无标签照片或请求人工反馈。由于它们的单位不同，我们分别保存预算：
\[
  (n_{\mathrm{lab}},n_{\mathrm{unlab}},n_{\mathrm{feedback}}).
\]
三个分量依次为人工标签数、无标签照片数和反馈轮次，只有给定单位成本后才可换算并相加。按模型疑问主动挑出的照片与均匀抽样照片，也可能在相同条数下提供不同信息，因而要记录查询策略和候选样本池。少样本任务中的支持集是供当前任务学习的带标签样本；比较其预算时，应同时固定类别数和每类样本数。总共五张照片不等于每个类别五张照片。

若目标标签完全不可用，预算再大的无标签样本也未必能够确认风险下降。考虑两个可能目标：照片分布相同，标签规则却相反。无标签程序收到的数据相同，正确预测应当不同，因此无法仅靠这些数据识别目标规则。输入距离和无标签预测一致性可以在额外假设下帮助筛选，但不能被当作已观察到的目标收益；这一信息边界可参考\BookChapterRef{chap:transfer-learning-theory}{迁移学习和持续学习的可学习性}。

\subsection{计算与存储收益的判断}
\label{sec:ket-reuse-resource-benefits}

达到同样准确率而少标注照片，仍可能需要更多计算。来源准备、候选搜索、特征提取和部署推理都要支付成本。为了比较整个服务期限，令$C_{\mathrm{prep}}$为一次性准备成本，$C_{\mathrm{sel},k}$与$C_{\mathrm{adapt},k}$为目标任务$k$的选择和适应成本，$q_k$为请求数，$c_{\mathrm{infer},k}$为一次推理成本。假定这些量已换算到同一计算或计费单位，先逐项加总，再将准备成本除以实际任务数：
\begin{equation}
  \begin{aligned}
    C_{\mathrm{reuse}}(m)
    &=C_{\mathrm{prep}}+
      \sum_{k=1}^{m}
      \bigl(C_{\mathrm{sel},k}+C_{\mathrm{adapt},k}
      +q_kc_{\mathrm{infer},k}\bigr),\\
    \overline C_{\mathrm{prep}}(m)&=\frac{C_{\mathrm{prep}}}{m}.
  \end{aligned}
  \label{eq:ket-reuse-amortization}
\end{equation}
式中$m$是核算期限内实际采用该方案的任务数。摊销只改变准备成本的分配，不改变总支出，也不证明复用比对照便宜。已经失败或被淘汰的试训同样消耗资源，应纳入相应选择与适应账本。不同型号设备上的GPU小时不能直接视为等价计算量；可以固定设备与计费条件，或分别报告设备时长、运算量和金额。

用同一种GPU构造一个成本比较。来源准备花120 GPU小时，每个目标的选择和适应分别花2与8 GPU小时；从头训练每个目标花30 GPU小时。规定两条路线达到相同效果、服务同样请求量，推理成本相同，在差值中相消。复用总成本为$120+10m$，从头训练为$30m$。要使准备投入不超过目标端累计节省，需要
\begin{equation}
  120+10m\leq30m
  \quad\Longleftrightarrow\quad m\geq6
  \label{eq:ket-reuse-break-even}
\end{equation}
成立。六个任务时两者均花180 GPU小时，这只是收支相等，严格节省要等到$m>6$。一个任务时复用花130、对照花30 GPU小时；十个任务时分别为220与300，复用少花80。若复用每个任务还增加5 GPU小时推理，比较式变为$120+15m\leq30m$，到八个任务才收支相等。由此可见，同一准备投入的资源结论取决于实际服务期限。

如果来源已经训练完成且允许直接使用，当前项目不需要重新支付那120 GPU小时。这属于增量决策：历史支出已经发生，当前要比较尚可改变的新增支出。下载、转换、验证和授权等新增成本仍须计入。从模型整个生命周期核算时，来源训练成本又必须保留。两种口径都可报告，但要注明边界，不能在复用侧排除准备投入、在对照侧却计入全部准备。

部署成本还包括不参与训练的权重。设共同特征提取网络即骨干有1亿参数，每个目标头有100万参数，每参数保存为4字节。采用十进制MB且忽略文件元数据与归档副本，十个全量微调模型共占$10\times404=4040$ MB；十个任务共用冻结骨干、分别保存头，则占$400+10\times4=440$ MB。共享减少权重副本，却没有免去每张新照片的骨干前向计算。为$n$张照片缓存$d$维特征、每元素$b$字节，还需$ndb$字节；训练时梯度、优化器状态和中间激活的峰值内存另计。多个来源同时推理时也须重新加总，不能用可训练参数量代替这些资源。

人工核验、标注、计算和存储各有不同单位，应先分别列示。给定设备单价、存储期限、请求量和人工价格后，才可换算货币支出；准确率若要进入效用，还须说明错误的应用代价。表~\ref{tab:ket-reuse-benefit-protocol}~汇总前面的三个比较问题，提醒我们固定什么、比较什么。同一来源在一种预算下有利，不会自动在另一种效果要求或服务期限下仍然有利。

\begin{table*}[htbp]
  \centering
  \small
  \begin{tabularx}{\linewidth}{lXXXXX}
    \toprule
    收益 & 固定条件 & 比较量 & 必要对照 & 准备成本 & 常见误判 \\
    \midrule
    预测效果 & 目标标签、结构及调参机会 & 同一终点的错误率差 & 来源初始化与随机初始化；冻结随机特征可作补充 & 来源训练及目标端选择分别记录 & 第10步领先18个百分点，误称第100步仍领先 \\
    样本效率 & 效果阈值、类别覆盖与标注粒度 & 达阈值的标签、无标签及反馈预算 & 从头训练的完整学习曲线 & 筛选30个标签不能漏记 & 把80个适应标签当作全部标注支出 \\
    计算存储 & 达成效果、设备、请求量和期限 & 总计算、峰值内存、权重与缓存 & 同效果的从头训练或其他可执行方案 & 来源120 GPU小时按实际任务数摊销 & 一次任务更贵，却借十次任务的平均值声称省算力 \\
    \bottomrule
  \end{tabularx}
  \caption{三种收益所需的比较条件。预测效果固定预算，样本效率固定效果，资源成本还固定服务期限与计费边界；表内数字来自本章教学构造。各行的比较量单位不同，不直接合成来源排名。}
  \label{tab:ket-reuse-benefit-protocol}
\end{table*}

\section{可复用知识的判断}
\label{sec:ket-reuse-compatibility}

明确收益口径以后，我们还需要判断哪些来源值得试用。物种模型见过目标地区的照片，未必保留了病害纹理；保留了纹理，也未必通过开放的物种预测接口提供这些信息。以下对同一候选分别检查输入输出是否相容、表示能否支持目标预测，以及任务之间是否共享可利用的结构。这三方面可以共同限制或支持复用，不是一组互斥来源类别。输入接近、标签语义一致与预测规则可共用，也分别需要证据。

\subsection{输入输出关系的判断}
\label{sec:ket-reuse-input-output}

先检查目标照片怎样进入来源模型。尺寸、通道和编码可被读取，只说明计算接口接得上，还要核对坐标所表达的观测含义。例如，普通RGB照片记录红、绿、蓝通道的光学信息，近红外图像记录另一波段的响应；即使两者都整理成三通道数组，相同位置的数值也未必对应同一观测。覆盖检查则考察哪些目标情形在来源中出现过。来源只含花朵正面，目标包含夜间或严重遮挡照片时，应具体登记这些缺口，不能由全局平均距离较小推断每个目标区域都有依据。

即使输入分布完全相同，还要检查输出要求。用四种等概率属性组合构造反例，令$\bm{x}=(x_1,x_2)\in\{-1,1\}^2$，$x_1$代表颜色属性，$x_2$代表斑点属性，两侧都均匀采样。来源预测$y_S=\mathbb{I}(x_1=1)$，目标预测$y_T=\mathbb{I}(x_2=1)$；指示函数$\mathbb I$在括号内条件成立时取1，否则取0。复制来源规则，四种目标组合中两种答错，错误率为$1/2$。若目标改为$y_T=1-y_S$，则四种全部答错。输入分布的任何合法距离都为零，却没有建立条件标签关系的相容性。

这个反例否定的是直接沿用来源预测规则，并没有否定全部来源表示。若中间层仍保留$(x_1,x_2)$，目标可以重新读取斑点信息；若只保留$x_1$，独立的$x_2$就无法从该表示恢复。相同输出名称也需要核对含义：来源把疑似病斑标为健康，目标规定疑似病斑属于异常并须送检时，即使两侧仍使用“健康／异常”标签，标签规则已经变化。选择参数起点、固定表示和旧预测头时，须分别解释这些变化的影响。

没有依据与已有依据冲突，还应分别记录。来源只见过甲、乙物种，目标加入丙时，我们缺少来源对丙的覆盖证据，但还不能断定其纹理或形态表示有害。如果一批已确认属于丙的照片在来源规范中被标为乙，则出现具体的标签冲突。前一种情形需要目标样本或未知类处理规则，后一种情形需要核对类别映射和标注标准；把二者都称为“不相似”会遗漏所需的不同操作。

输入输出检查的结果可以写成具体记录：哪些目标情形在现有来源样本中未见，哪些标签映射已由规范或配对样本确认，哪些观测对应仍待验证。例如，夜间样本暂无覆盖，甲、乙定义保持一致，丙不在来源输出集合内。这些记录分别指向补充数据、沿用标签和扩展输出的需要。图~\ref{fig:ket-reuse-coverage-relations}~左侧示意覆盖关系；区域重叠只说明输入可能共同出现，不能证明标签机制一致。

覆盖和共同低风险预测器的存在都是需要证据支持的条件。例如，在共同标签机制和已知且有界的密度比下，源样本覆盖目标区域才支持相应的加权风险估计。本节只形成适用条件及证据缺口清单，尚未实施适应。分布校正与跨观测对应交由第~\ref{chap:ket-update-maintenance}~章解释，新类别和新输出的语义要求交由第~\ref{chap:ket-transfer-adaptation}~章解释。

\subsection{表示与预测关系的判断}
\label{sec:ket-reuse-representation}

现在把复用对象固定为来源网络的一层表示。我们需要知道病害信息是否仍在，以及能否被计划采用的预测头读出。可以冻结第$\ell$层特征提取器$\phi_\ell$，仅用目标标签训练容量受控的新头，这称为\term{表示探测}（representation probing）。探测检验的是某种读取方式下的信息可用性；它不枚举所有可能预测器，因而不能由一个头失败就断定表示丢失了全部目标信息。

以$c$类目标为例，把$n$个带标签训练输入变成$\bm z_i=\phi_\ell(\bm x_i)\in\mathbb R^{d_\ell}$，并把标签编号为$y_i\in\{1,\ldots,c\}$。限定预测头的类别评分为$\bm W\bm z_i+\bm b$，再用softmax把评分转成归一化概率。为使正确类别获得高概率，最小化其负对数概率的平均；加入权重平方惩罚，限制头依靠过大系数拟合少量标签：
\begin{equation}
  \begin{aligned}
    (\hat{\bm W},\hat{\bm b})
    =\argmin_{\bm W,\bm b}\Bigl\{
    &-\frac{1}{n}\sum_{i=1}^{n}
    \log\bigl[\operatorname{softmax}
    (\bm W\bm z_i+\bm b)\bigr]_{y_i}\\
    &+\lambda\lVert\bm W\rVert_F^2\Bigr\},
  \end{aligned}
  \label{eq:ket-reuse-linear-probe}
\end{equation}
候选范围为$\bm W\in\mathbb R^{c\times d_\ell}$和$\bm b\in\mathbb R^c$；$\lVert\bm W\rVert_F^2$是全部权重元素的平方和，$\lambda>0$控制惩罚强度。固定特征后，这个目标对头参数是凸的，可用梯度法求解，但停止时仍须检查优化误差。若训练集涵盖所有类别，正则化权重后可取得有限最优头；所有类别评分同时加一个常数不改变概率，可固定$b_c=0$消除这种偏置平移。类别缺失或取消正则时须另查最优值是否取到。正则强度与停止规则通过不参与头拟合的验证数据选择，来源参数始终不变。新输入经过同样预处理、$\phi_\ell$和新头得到概率，再选最大概率类别；缓存特征时也须冻结归一化统计量，保持各候选的提取规则一致。

探测低分可能来自信息丢失，也可能来自线性头读不出仍在的信息、标签不足或优化未完成。改用更复杂的头可以检查读取限制，却同时增加了目标学习本身的能力，不能继续把成绩全部归给来源表示。不同层的比较因此须记录特征维数、空间特征如何汇总成向量、头参数量、标签量和优化预算，并使用相同目标划分。只有这些条件明确，探测结果才有可比较的含义。

前面的属性反例还没有比较不同表示，现在补上这一点。$\phi_A(\bm x)=(x_1,x_2)$与$\phi_B(\bm x)=x_1$都能支持来源规则$y_S=\mathbb I(x_1=1)$，来源准确率同为100\%。对目标$y_T=\mathbb I(x_2=1)$，A可以按第二个坐标的符号正确分类。B对固定$x_1$合并了两种$x_2$，而对应的两个目标标签等概率出现，任何只读B的预测器最高准确率都是50\%。这里有明确的信息合并证据，不是仅从一个未训练好的头猜测原因，也不能由相同来源成绩预见这种差别。

模块接回网络后还可能出现另一类困难。Yosinski等迁移不同层，比较冻结、微调和同任务重新连接的对照，以区分来源任务特异性与切断相邻层协同计算造成的优化问题。因此，复制模块后效果下降时，还应检查下游网络能否重新利用它，而不只检查模块所含信息。该结论来自论文中的卷积网络和任务划分，不能据此指定所有架构中必然通用的一层\scite{yosinski2014}

能够实施哪种探测取决于来源接口。只有黑盒预测时，可以核对类别语义、概率归一化和预测置信度与目标标签的关系，却没有中间层线性可分性的观测证据。预测头没有提供病害信息，也不证明早期表示没有该信息。取得特征则可以做表示探测，但仍不等于获得修改源参数的权限。第~\ref{sec:ket-reuse-fixed-score}~节将分别讨论概率接口和特征接口的评分；这些评分最终仍要接受目标对照的检验。

\subsection{任务结构关系的判断}
\label{sec:ket-reuse-task-relations}

对于已有的物种、病害和定位模型，还可以检查它们是否提供目标所需的子问题结构。物种识别比较形态，病害判断寻找局部异常，定位输出花朵的空间位置；输入同为照片，输出要求不同。定位信息若能被当前程序接入物种识别，就可能帮助把花朵与背景分开，但这种作用还需要试验。定位任务不必成为最终交付目标，也可以作为辅助来源，条件是它的结果或表示确实进入目标计算。

要检验任务$a$是否帮助任务$b$，我们固定目标$b$的损失与预算，把从$a$的知识出发得到的$f_{a\to b}$和无迁移预测器$f_{0,b}$比较。沿用风险下降为正的方向，定义有向收益
\begin{equation}
  G_{a\to b}=R_b(f_{0,b})-R_b(f_{a\to b}).
  \label{eq:ket-reuse-directed-gain}
\end{equation}
这里$R_b$表示在目标$b$分布上按其损失计算的风险，$G_{a\to b}>0$才表示在该协议下有帮助。反向试验把目标、损失、对照与来源信息都换了，因此一般没有$G_{a\to b}=G_{b\to a}$的依据。不同目标的损失单位不一致时，还不能直接相加收益，或用同一箭头粗细比较大小。

图~\ref{fig:ket-reuse-coverage-relations}~右侧用固定预算构造三个方向。病害目标的100张验证照片中，从头训练错20张，物种来源适应后错18张，故物种$\to$病害收益为2个百分点。物种目标另有100张验证照片，从头训练错10张，病害来源适应后错11张，反向收益为负1个百分点；定位来源适应后错6张，收益为4个百分点。两侧目标均按零一损失计分，但使用各自对照。未测方向不画箭头，不视为收益为零。

\begin{figure*}[htbp]
  \centering
  % Teaching construction, not a measured distribution or task graph.
% Coverage rectangles encode set membership only, not probability mass.
% Directed gains (percentage points):
% species -> disease: (20 - 18) / 100 = +2;
% disease -> species: (10 - 11) / 100 = -1;
% localization -> species: (10 - 6) / 100 = +4.
% All other directions are unmeasured, not zero.
\begin{tikzpicture}[
  x=1cm,y=1cm,font=\figurefont\small,
  task/.style={
    rectangle,draw=BookInk,fill=BookPaper,line width=0.8pt,
    minimum width=24mm,minimum height=9mm,align=center,inner sep=2mm
  },
  relation/.style={
    draw=BookInk,line width=0.9pt,
    -{Latex[length=2mm]},rounded corners=1.2mm
  },
  edge label/.style={font=\figurefont\footnotesize,text=BookInk,inner sep=1mm}
]
  \node[anchor=west,font=\figurefont\small\bfseries,text=BookInk]
    at (0.1,4.65) {A\quad 预测输入覆盖};
  \node[anchor=west,font=\figurefont\small\bfseries,text=BookInk]
    at (8.1,4.65) {B\quad 有向任务收益};
  \fill[FigureInput!10] (0.5,0.8) rectangle (4.1,3.1);
  \fill[FigureOutput!10] (2.8,1.4) rectangle (4.1,3.1);
  \fill[BookAmber!12] (4.1,1.4) rectangle (6.8,3.6);
  \fill[BookAmber!12] (2.8,3.1) rectangle (4.1,3.6);
  \draw[BookBlue,line width=0.9pt] (0.5,0.8) rectangle (4.1,3.1);
  \draw[BookTeal,dashed,line width=0.9pt] (2.8,1.4) rectangle (6.8,3.6);
  \node[text=BookBlue] at (1.65,2.1) {来源};
  \node[text=BookTeal] at (5.55,3.2) {目标};
  \node[align=center,font=\figurefont\footnotesize,text=BookInk]
    at (3.45,2.2) {共同\\覆盖};
  \node[align=center,font=\figurefont\footnotesize,text=BookAmber]
    at (5.5,2.1) {目标中的\\来源缺口};
  \node[anchor=west,text=BookMuted,font=\figurefont\footnotesize]
    at (0.5,0.2) {重叠不代表标签机制相同};

  \node[task] (species) at (9.45,2.4) {物种识别};
  \node[task] (disease) at (14.7,2.4) {病害判断};
  \node[task] (localization) at (9.45,0.3) {花朵定位};
  \draw[relation] (species.north)
    -- (9.45,3.85) -- (14.7,3.85) -- (disease.north);
  \node[edge label,anchor=south] at (12.075,3.9) {$+2$};
  \draw[relation] (disease.west) -- (species.east);
  \node[edge label,anchor=south] at (12.075,2.45) {$-1$};
  \draw[relation] (localization.north) -- (species.south);
  \node[edge label,anchor=west] at (9.55,1.35) {$+4$};
  \node[align=left,anchor=west,text=BookMuted,font=\figurefont\footnotesize]
    at (11.45,0.3) {单位：百分点\\未画方向：未测};
\end{tikzpicture}

  \caption{输入覆盖和有向任务收益的教学示意。左侧实线框为来源、虚线框为目标，面积不是概率，重叠不是标签规则一致的证据。右侧箭头从来源指向目标，数字为相对该目标无迁移对照的准确率增量，单位为百分点；未画出的方向尚未测量。}
  \label{fig:ket-reuse-coverage-relations}
\end{figure*}

图中的试验先固定来源学习结果，再适应目标。若改为两个任务共同训练，双方会继续改变共享表示，也可能竞争容量或提出冲突更新要求，评价对象已经不同。因此，现成定位器帮助物种分类，不意味着在物种训练中加入定位损失就有相同收益。第~\ref{sec:ket-sharing-affinity}~节讨论共同学习时任务之间的影响；这里的有向箭头只对已规定的顺序迁移协议成立。

上述箭头比较具体任务；如果要复用前章介绍的参数先验或学习规则，还需问历史任务是否代表将来的学习问题。沿用第~\ref{sec:ket-overview-learning-settings}~节，$\mathcal T_k$规定输入输出语义、任务内分布$\mathcal P_k$和损失$\ell_k$，$\Pi_{\mathrm{hist}}$规定怎样抽取历史任务。再固定支持集与反馈协议，使规则$U$根据$S_k$和既有状态产生预测器，简写为$U(S_k)$。要衡量规则在历史任务范围中的平均风险，就先在每项任务内计算$R_k$，再对任务抽样取期望：
\[
  \mathbb E_{\mathcal T_k\sim\Pi_{\mathrm{hist}}}
  \mathbb E_{S_k\mid\mathcal T_k}[R_k(U(S_k))]
\]
式中$R_k$为任务$\mathcal T_k$的风险，支持集抽样的随机性也按已定协议包含在平均中。跨任务平均还要求各损失单位可比，或先规定归一化与权重。这个量只描述$\Pi_{\mathrm{hist}}$所覆盖的任务。若历史任务都是每类数张照片的物种分类，新任务要求病斑分割，支持集由类别标签变成区域标注，输出与损失也变化，原平均不能支持对此新任务的快速适应承诺。

检查规则的适用性时，应核对历史任务包含了哪些类别重组、光照变化和支持集规模，而非只数历史任务个数。在固定类别规则下反复抽样，没有自动产生适应新输出结构的证据。目标变化不在历史范围内时，应安排针对该目标的试训，或明确保留不确定性。这一检查也为第~\ref{sec:ket-transfer-task-splits}~节的任务划分和学习策略准备提供依据。

\section{可复用知识的选择}
\label{sec:ket-reuse-selection}

完成适用性检查后，我们得到候选对象和各自的证据缺口，还没有得到可靠排名。把每个候选都完整适应一次可能耗尽预算，廉价评分又可能漏掉后期才见效或需要组合的来源。选择因此包括三项相连工作：估计候选价值，在可行范围内选择来源与组合，再用新证据检验并调整决定。候选可以是整个模型、某层表示与指定预测头，或按规则纳入的样本子集；估计时必须评价计划真正复用的对象。

\subsection{候选知识价值的估计}
\label{sec:ket-reuse-value-estimation}

假如当前目标是提高给定标签预算下的终点准确率，估计就应尽量对应这一效果，而非来源成绩或早期下降速度。估计达90\%所需标签、估计整个部署成本，则需要另一种观测和核算。下面先讨论目标效果的证据，再将预算作为选择约束。固定来源模型的评分不运行实际适应过程，成本较低但仍需标签、前向计算或数值求解；目标试训花更多计算观察适应过程，二者的证据范围不同。

\subsubsection{基于固定模型的价值估计}
\label{sec:ket-reuse-fixed-score}

只开放来源类别概率时，我们无法探测中间表示，但仍可检查这些概率是否区分目标类别。例如，来源预测物种，目标预测病害，不要求两组类别同名，却可以观察某些来源响应是否经常对应同一目标标签。\term{对数期望经验预测}（log expected empirical prediction，LEEP）据此建立经验对应，再为对应后的目标预测计算分数。来源模型保持固定；方法需要完整的离散类别概率和带标签的目标分类样本\scite{ket-nguyen2020-leep}

固定非空评分集$S_{\mathrm{score}}=\{(\bm x_i,y_i)\}_{i=1}^{n}$，令$Y_T$和$Y_S$分别为目标、来源类别集合。来源对输入$\bm x_i$给出$q_s(\bm x_i)\geq0$，且$\sum_{s\in Y_S}q_s(\bm x_i)=1$。为估计“来源响应$s$对应哪个目标$y$”，把目标标签为$y$的样本贡献的$s$概率相加并除以$n$，再按来源列归一化：
\begin{equation}
  \begin{aligned}
    \hat p(y,s)&=\frac1n\sum_{i:y_i=y}q_s(\bm x_i),\\
    \hat p(s)&=\sum_{y\in Y_T}\hat p(y,s),\qquad
    \hat p(y\mid s)=\frac{\hat p(y,s)}{\hat p(s)}.
  \end{aligned}
  \label{eq:ket-reuse-leep-joint}
\end{equation}
每个样本把质量$1/n$分配给不同来源类别，总质量不变，所以$\hat p(y,s)$是归一化经验联合分布。$\hat p(s)$是第$s$列的边缘质量，$\hat p(y\mid s)$描述该列中的目标标签比例；它们不是来源训练集上的真实联合分布。若某列$\hat p(s)=0$，非负性保证评分集中该列的$q_s$全为零，可在计算评分时跳过，不执行除零。在新输入上使用EEP还需约定零质量来源列的目标分布，不能把评分集外的概率质量直接丢弃。

得到对应表后，对每张照片按来源概率加权各列的目标分布，形成\term{期望经验预测器}（expected empirical predictor，EEP），即把来源响应转成目标预测的混合规则。用它给正确标签分配的概率计算平均对数分数：
\begin{equation}
  \begin{aligned}
    p_{\mathrm{EEP}}(y\mid\bm x)
    &=\sum_{s\in Y_S}\hat p(y\mid s)q_s(\bm x),\\
    s_{\mathrm{LEEP}}
    &=\frac1n\sum_{i=1}^{n}
      \log p_{\mathrm{EEP}}(y_i\mid\bm x_i).
  \end{aligned}
  \label{eq:ket-reuse-leep-score}
\end{equation}
以下$\log$均为自然对数。概率不超过1，因此LEEP不大于零，数值越大表示在这批评分样本上正确标签得到的概率越高；零要求每个正确标签概率都是1。选择对数概率是在考察目标标签的预测似然，并非直接计算准确率。条件表和分数使用同一批标签，故LEEP是候选价值的代理量，不是独立测试结果。

表~\ref{tab:ket-reuse-score-data}~构造四张病害评分照片，局部将目标标签编号为0和1。候选A、B都提供两个来源类别的概率和一维冻结特征。为保证两种接口相容，规定A的头为$q_1=\operatorname{sigmoid}(-\log(9)z_A)$，其中$\operatorname{sigmoid}(u)=1/(1+\exp(-u))$；B的头恒输出$(0.5,0.5)$。两者均令$q_2=1-q_1$。后文将用同一组样本分别考察来源概率和冻结特征。

\begin{table}[htbp]
  \centering
  \small
  \begin{tabularx}{\linewidth}{rrXXrr}
    \toprule
    $i$ & $y_i$ & A的$(q_1,q_2)$ & B的$(q_1,q_2)$ & $z_{A,i}$ & $z_{B,i}$ \\
    \midrule
    1 & 0 & $(0.9,0.1)$ & $(0.5,0.5)$ & $-1$ & $-1$ \\
    2 & 0 & $(0.9,0.1)$ & $(0.5,0.5)$ & $-1$ & $1$ \\
    3 & 1 & $(0.1,0.9)$ & $(0.5,0.5)$ & $1$ & $-1$ \\
    4 & 1 & $(0.1,0.9)$ & $(0.5,0.5)$ & $1$ & $1$ \\
    \bottomrule
  \end{tabularx}
  \caption{同一目标评分集的概率接口与特征接口。$y=0$为健康、$y=1$为病害；A、B的来源概率均按行归一化，末两列为各自的一维冻结特征。所有数值为教学构造，供LEEP与LogME复算，不是真实图像测量。}
  \label{tab:ket-reuse-score-data}
\end{table}

按目标类别0、1排列行，按来源类别1、2排列列，将A的概率代入经验对应公式，得到
\[
  \bigl[\hat p_A(y,s)\bigr]=
  \begin{pmatrix}0.45&0.05\\0.05&0.45\end{pmatrix},
  \qquad
  \bigl[\hat p_A(y\mid s)\bigr]=
  \begin{pmatrix}0.9&0.1\\0.1&0.9\end{pmatrix}.
\]
两列边缘质量均为0.5，所以列归一化把第一张表的元素乘以2。第一张照片的正确标签0得到概率$0.9\times0.9+0.1\times0.1=0.82$，其余三张也相同，故$s_{\mathrm{LEEP},A}=\log0.82\approx-0.198451$。B的联合表每格0.25，条件表每格0.5，分数为$\log0.5\approx-0.693147$。在同一评分集上，LEEP把A排在B之前，依据是A的来源响应与病害标签有更明确的经验对应。

B的恒定输出说明一个边界：EEP只能恢复评分集中的类别频率。此时LEEP为$\sum_y\hat p(y)\log\hat p(y)$，即负的经验标签熵；若绝大多数照片健康，这个数也会接近零，却没有识别稀有病害的能力。来源应在同一目标集上比较，不能按不同任务的原始分数作通用排名。若接口仅返回最大概率类别，可以将硬输出编码成独热向量另作计算，但概率信息已经丢失，不等价于完整概率接口上的原始LEEP。

LEEP与训练似然的关系还依赖候选范围。如果固定表示上的目标预测头集合包含EEP，那么该集合可达到的训练平均对数似然至少等于LEEP，因为可以直接选EEP。这一比较针对同一评分样本的似然，不是独立测试准确率下界。普通线性头未必包含这个混合预测器，数值求解也未必达到集合中的最佳值；解冻后表示改变，更需要重新检验。

若候选仅提供编码器而没有分类头，LEEP缺少来源类别概率；目标为位置坐标等连续量时，离散对应也不直接适用。可以转而问冻结特征能否在线性预测模型中解释目标标签。\term{最大证据的对数}（logarithm of maximum evidence，LogME）为线性权重指定先验，对不同权重的标签似然加权积分，得到模型证据，再选择证据支持的超参数。积分以参数分布保留不确定性，相应参数学习原理可参考\BookChapterRef{chap:pgm-parameter-learning}{参数学习}；这里给出评分所需的模型与积分\scite{ket-you2021-logme}

从一个候选取得$n$个$d$维特征，按样本排列成$\bm Z\in\mathbb R^{n\times d}$。先用一维实值标签$\bm y\in\mathbb R^n$说明计算。我们以$\bm Z\bm w$预测标签，用独立高斯噪声描述剩余偏差，并用零均值高斯先验限制权重尺度：
\begin{equation}
  \begin{aligned}
    \bm w&\sim\mathcal N(\bm0,\alpha^{-1}\bm I_d),\\
    \bm y\mid\bm Z,\bm w
    &\sim\mathcal N(\bm Z\bm w,\beta^{-1}\bm I_n),
    \qquad \alpha,\beta>0.
  \end{aligned}
  \label{eq:ket-reuse-logme-model}
\end{equation}
其中$\bm w\in\mathbb R^d$，$\bm I_d$与$\bm I_n$分别为相应维数的单位矩阵；$\alpha,\beta>0$是方差的倒数，分别控制先验权重与观测噪声的精度。先验偏好小权重，并未固定$\bm w=\bm0$。原式不含显式截距，因此中心化、标准化或增加常数特征都会改变模型，应对候选采用固定并可复现的处理规则。

只比较某个最佳权重的训练似然，没有计入先验允许的其他权重是否也能解释标签。模型证据改为按先验对整个权重空间$\mathbb R^d$加权积分，计算该模型对观测标签分配的密度：
\begin{equation}
  p(\bm y\mid\bm Z,\alpha,\beta)
  =\int p(\bm y\mid\bm Z,\bm w,\beta)
       p(\bm w\mid\alpha)\,\mathrm d\bm w.
  \label{eq:ket-reuse-logme-evidence}
\end{equation}
由于先验与噪声都是高斯分布，这个积分可以配方求出。令$\bm A=\alpha\bm I_d+\beta\bm Z^\top\bm Z\in\mathbb R^{d\times d}$，它在$\alpha>0$时正定；后验均值为$\bm m=\beta\bm A^{-1}\bm Z^\top\bm y\in\mathbb R^d$。合并两个高斯指数中的平方项，得到
\[
  \begin{aligned}
    &\alpha\lVert\bm w\rVert_2^2+
      \beta\lVert\bm y-\bm Z\bm w\rVert_2^2\\
    &\quad=(\bm w-\bm m)^\top\bm A(\bm w-\bm m)
       +\alpha\lVert\bm m\rVert_2^2
       +\beta\lVert\bm y-\bm Z\bm m\rVert_2^2.
  \end{aligned}
\]
第一项是以$\bm m$为中心的权重二次型，其指数因子积分为$(2\pi)^{d/2}|\bm A|^{-1/2}$。余下两项不含积分变量$\bm w$，可以提出积分号。连同先验与似然的归一化常数取对数，得到
\begin{equation}
  \begin{aligned}
    L(\alpha,\beta)
    &=\frac n2\log\beta+\frac d2\log\alpha
       -\frac n2\log(2\pi)-\frac12\log|\bm A|\\
    &\quad-\frac\beta2\lVert\bm y-\bm Z\bm m\rVert_2^2
       -\frac\alpha2\lVert\bm m\rVert_2^2.
  \end{aligned}
  \label{eq:ket-reuse-logme-log-evidence}
\end{equation}
这里$|\bm A|$为行列式，$\bm y-\bm Z\bm m$为后验均值预测留下的残差。残差项反映拟合，权重项与行列式项反映先验尺度和被数据约束的权重范围，所以证据不只按训练误差排序。这一处理仍依赖线性预测、高斯噪声和各向同性先验，即先验在各权重方向使用相同方差；它不是任意模型都不会过拟合的保证。

候选特征固定后，还需确定先验与噪声尺度。目标是最大化对数证据$L(\alpha,\beta)$，不是进一步训练来源网络。对证据求偏导并令其为零，可整理出下列固定点关系。令$\lambda_j$为$\bm Z^\top\bm Z$的非负特征值，第$t$步先用当前$\alpha^{(t)},\beta^{(t)}$计算后验均值$\bm m^{(t)}$，再计算各特征方向合计的有效维数$\gamma^{(t)}$并更新：
\begin{equation}
  \begin{aligned}
    \gamma^{(t)}&=\sum_{j=1}^{d}
       \frac{\beta^{(t)}\lambda_j}{\alpha^{(t)}+\beta^{(t)}\lambda_j},\\
    \alpha^{(t+1)}&=
       \frac{\gamma^{(t)}}{\lVert\bm m^{(t)}\rVert_2^2},\\
    \beta^{(t+1)}&=
       \frac{n-\gamma^{(t)}}{\lVert\bm y-\bm Z\bm m^{(t)}\rVert_2^2}.
  \end{aligned}
  \label{eq:ket-reuse-logme-update}
\end{equation}
若$\beta^{(t)}\lambda_j$远大于$\alpha^{(t)}$，该方向主要受数据约束，接近计为一维；反之主要受先验限制，接近计为零维。这解释了$\gamma^{(t)}$的作用。原算法从$\alpha=\beta=1$开始，复用特征分解计算这些量。实际求解应规定迭代上限与相对变化停止条件；零权重范数、零残差或边界最优须另外处理，不能直接除零。固定点成立只表示满足相应驻点关系，有限次更新并不保证取得全局最大证据。

分类标签不是连续观测，但LogME仍可把每类是否出现编码为回归目标。对$c$类构造$\bm Y\in\mathbb R^{n\times c}$，每行在所属类别位置取1，其余取0；对每列$\bm Y_{:,j}$分别使用上述模型、选择$\alpha_j,\beta_j>0$。先按$n$归一化各列对数证据，再按$c$平均，得到
\begin{equation}
  s_{\mathrm{LogME}}
  =\frac1c\sum_{j=1}^{c}\frac1n
    \sup_{\alpha_j,\beta_j>0}
    L_j(\alpha_j,\beta_j).
  \label{eq:ket-reuse-logme-score}
\end{equation}
这里$L_j$是第$j$列的对数证据，多输出回归也可逐列计算。独热标签的高斯回归只是分类代理，不是softmax分类器的精确边际似然。分数越大表示标签在该代理模型下得到的密度越高；连续密度可以大于1，其对数不受LEEP的零上界限制。采用上确界允许超参数趋向边界而没有有限最大点；某些退化数据甚至可能使证据无界，数值实现必须说明有限搜索范围及边界处理，不能把无界值当作可靠排名。比较候选时应固定目标集合、标签单位、预处理和搜索规则。

再用表~\ref{tab:ket-reuse-score-data}~计算特征评分，这次不读取来源头。病害列为$\bm y=(0,0,1,1)^\top$，A的特征列为$\bm z_A=(-1,-1,1,1)^\top$，故$\bm z_A^\top\bm z_A=4$、$\bm z_A^\top\bm y=2$、$\bm y^\top\bm y=2$。积分后的标签密度仍为零均值高斯，其协方差是噪声协方差加上随机线性预测的协方差：
\[
  \bm C=\beta^{-1}\bm I_4+\alpha^{-1}\bm z_A\bm z_A^\top.
\]
$\bm z_A$的长度为2，沿单位向量$\bm u=\bm z_A/2$的方差为$v=\beta^{-1}+4/\alpha$；三个正交方向的方差均为$r=\beta^{-1}$。标签在$\bm u$上的平方投影为1，余下平方长度也为1。因此行列式为$vr^3$，二次型为$1/v+1/r$，对数密度可化为
\[
  L_A=-\frac12\left[
  4\log(2\pi)+\log v+3\log r+\frac1v+\frac1r\right].
\]
最大化$L_A$等价于分别最小化$\log v+1/v$和$3\log r+1/r$，导数分别在$v=1$与$r=1/3$变号。这两个值满足$v>r>0$，因此可对应$\beta=3$、$\alpha=6$。后验均值为标量$1/3$，残差平方和为$10/9$，代入对数证据得到
\[
  \frac{L_A}{4}
  =-\frac{4\log(2\pi)-3\log3+4}{8}
  \approx-1.006959.
\]
健康列$(1,1,0,0)^\top$的投影符号相反，但两个平方长度相同，故对数证据相同，上式也是A的两类平均LogME。从$(1,1)$出发，固定点更新前两次得到$(5,3.076923)$和$(5.625,3.035552)$。它们在本例中向解析最优点$(6,3)$靠近；这段可复算迭代不代替一般收敛证明。

B的$\bm z_B=(-1,1,-1,1)^\top$与两列标签均正交。在该方向增加方差，不减少标签二次型，却增加行列式，所以最佳证据在$\alpha\to\infty$时逼近，相当于不使用该特征。此时协方差只剩$r\bm I_4$；标签平方长度为2，最优$r=2/4=1/2$，即$\beta=2$，于是
\[
  s_{\mathrm{LogME},B}
  =-\frac12\bigl(\log\pi+1\bigr)
  \approx-1.072365.
\]
B没有有限的最优$\alpha$，报告应注明边界解；有限搜索中很大的$\alpha$只是近似。LogME仍把A排在B之前，但它与LEEP使用不同预测模型和计分尺度。不能将A的LEEP值$-0.198451$减去其LogME值$-1.006959$来判断哪个评分方法更好；方法优劣须由它们对实际迁移效果的预测能力比较。

接口也限制哪些候选可比较。B若仅开放预测，LEEP可算而LogME缺少特征；仅提供无分类头编码器的候选，则缺少原始LEEP所需的类别概率。取得来源前向结果后，LEEP主要累计来源概率与目标标签对应，LogME还需存储或分批处理特征并做分解。直接构造并分解$d\times d$矩阵$\bm Z^\top\bm Z$约需$O(nd^2+d^3)$运算，多列标签可复用分解；特征维数很大时还需采用秩结构或其他求解路径。来源前向、数据存储和数值求解都没有因“不微调”而消失。

检验评分是否有用，需要把代理排名与计划使用方式下的实际目标效果比较。假设代理顺序为A、B、C，充分适应后的顺序为B、A、C，没有并列。三对比较中两对一致、一对相反，Kendall秩相关即一致对数减相反对数再除以总对数，为$1/3$；排名有正相关，首选却仍错了。即使所有次序都正确，全部来源也可能输给无迁移对照。候选范围、评分标签量、排序质量和实际风险收益必须分别报告。

\subsubsection{基于目标试训的价值估计}
\label{sec:ket-reuse-target-trials}

固定评分没有执行计划中的适应，因而无法观察解冻后的表示变化或实际优化困难。如果这些因素决定最终选择，可以用一部分预算运行目标学习程序，这称为\term{目标试训}（target-side trial training）。试训直接观察候选在指定使用方式下的结果，须先固定更新对象、预算和选择依据，不能随意截取各自最有利的一段曲线比较。

试训深度改变可观察的因素。冻结特征只训练轻量头，仍主要检查读取能力；允许局部参数更新，才能观察局部调整是否修正失配；运行完整适应，则可观察所有允许参数的共同变化和长程动态，也要支付对应成本。线性探测低分没有否定解冻潜力，冻结时高分也没有保证全量微调会保留有效结构。每个结果都须绑定其更新范围，而不只绑定来源名称。

把目标开发数据分为拟合集$S_{\mathrm{fit}}$和选择集$S_{\mathrm{select}}$，来源适应只用前者，候选决策用后者。尽量固定照片子集、类别构成、预处理、批量顺序、训练步数和查看频率，并事先规定初始化、优化器、学习率及正则的搜索规则。同一学习率便于隔离因素，却可能不适合所有候选；若比较各自适配后的能力，可在同样规模的预设网格中搜索。模型规模不同导致每步成本不同时，再限制设备时间或运算量，不把相同步数当成相同资源。

用表~\ref{tab:ket-reuse-score-data}~说明一次实际头更新，与前面的评分区分开来。这四张照片现在作为拟合集，另取四张不参与拟合的新照片作选择集；教学上规定新照片恰好具有同样的特征和标签组合。将二元标签编码为$t_i=2y_i-1$，固定预测评分为无截距标量头$wz_i$。用平方差衡量预测评分与标签编码的偏离，求平均并乘$1/2$以简化导数，按全批量梯度更新：
\begin{equation}
  H(w)=\frac{1}{2n}\sum_{i=1}^{n}(wz_i-t_i)^2,
  \qquad
  w^{(t+1)}=w^{(t)}-\eta H'(w^{(t)}).
  \label{eq:ket-reuse-trial-update}
\end{equation}
式中$n=4$，$w\in\mathbb R$为待拟合权重，$\eta$为学习率；$t_i$标记标签编码，上标$(t)$标记迭代。固定$w^{(0)}=0$、$\eta=1/2$，不加正则，仅训练两步。两候选的$n^{-1}\sum_i z_i^2$均为1，二次目标曲率为1，当前步长在稳定区间$(0,2)$内。A满足$n^{-1}\sum_i z_it_i=1$，故$H'_A(w)=w-1$，得到$w_A^{(1)}=1/2$、$w_A^{(2)}=3/4$；选择损失分别为$1/8$、$1/32$。B的特征与标签编码内积为零，$H'_B(w)=w$，两步权重仍为0，选择损失均为$1/2$。在这一头和两步预算下，A损失更低。

把评分转成类别时采用$\mathbb I(wz\geq0)$，平局归为1。A两步都正确分类这四张选择照片，准确率100\%；B恒预测1，准确率50\%。A的平方损失仍在下降，分类准确率却已经不变，说明两个指标不同。试训没有更新特征，结论只覆盖当前头、数据与预算；若允许更新B的编码器，上面的固定特征计算不再描述新过程。

若继续比较长程适应，应保存候选$j$在第$t$步的$\hat R_{\mathrm{select}}(f_j^{(t)})$和重复运行波动，不只保存最低点。另构造一组允许更新来源参数的结果：100张选择照片上，A在第2、8、16步错20、18、18张，B错22、16、14张。第二步A准确率80\%、B为78\%，但第16步A为82\%、B为86\%。前面已经区分早期与终点收益，这里新增的后果是：按第二步排名淘汰B，会使搜索根本看不到较好的终点。

可以用多阶段预算减少完整试训次数，但必须把淘汰成本也算进去。构造如下流程：8个候选各花1单位固定评分，留下4个；每个试训2步，留下2个；这2个从检查点继续到8步，再继续到16步。假定所有候选每步恰花1单位，保存检查点没有额外成本，总支出为
\[
  8\times1+4\times2+2\times(8-2)+2\times(16-8)
  =44.
\]
四项依次为初始评分、四个候选的前两步、两个候选新增六步及新增八步。作为对照，若跳过固定评分、直接把8个候选各试训16步，共花128单位。两条路线按各自实际执行的步骤计费；如果阶段间不保存检查点而重训，应计入全部重训步骤，不能仍减去已经执行的步数。

节省预算伴随漏选风险，尤其是只保留很少候选时。可为不同接口或表示类型、估计不确定性大的候选预留探索机会。若同一预定进度上各候选风险区间$[\underline r_j,\overline r_j]$同时覆盖相应风险，令$\delta\geq0$为允许差距，$\underline r_j>\min_h\overline r_h+\delta$表示候选$j$的最好可能风险仍比某个保留候选的最坏可能风险高出$\delta$，可据此淘汰当前进度上的劣势者，$h$遍历当前候选。自适应查看、重复比较或提前停止时，普通单次区间不自动维持整个过程的错误控制；应使用对允许查看时刻共同有效的区间，或固定阶段并校正多重比较。当前步有效也不意味着终点排序不变。没有学习曲线外推依据时，淘汰仍是承担误排风险的预算策略。

上述曲线与排名都利用了$S_{\mathrm{select}}$，它因此承担开发决策，不能再声称是独立最终测试。确认方案后，先锁定来源、更新范围、预算和停止规则，再用没有参与这些决定的数据评价。第~\ref{chap:ket-evaluation}~章进一步解释反复选择如何使开发成绩偏向被选中的方案。

\subsection{候选知识的选择与组合}
\label{sec:ket-reuse-choice}

候选价值已有估计以后，仍要核对使用所需的接口、授权和资源。一个高分来源可能不能在离线设备上运行，也可能需要同时保留其他来源才有收益。下面按选择单个来源还是来源集合展开，先规定硬约束，再比较可行方案。所选对象始终包含来源及其使用方式，不是脱离目标的模型排名。

\subsubsection{单来源知识的选择}
\label{sec:ket-reuse-single-source}

先考虑只采用一个来源的离线部署。把实际可取得的候选编号组成有限集合$J$，加入$j=0$代表无迁移方案，得到$J_0=J\cup\{0\}$。对每个方案估计同一终点风险$\hat r_j$，并记录选择后尚需支付的适应成本$c_j$、部署存储$m_j$和接口及授权可用标志$a_j\in\{0,1\}$。无迁移方案也按同一边界记账，其训练成本不设为零。给定剩余计算预算$B_c$和存储预算$B_m$，只在满足约束的候选中选预计风险最小者：
\begin{equation}
  \begin{aligned}
    j^\star&\in\argmin_{j\in J_0}\hat r_j\\
    \text{满足}\quad
    &c_j\leq B_c,\quad m_j\leq B_m,\quad a_j=1.
  \end{aligned}
  \label{eq:ket-reuse-single-choice}
\end{equation}
取最小风险而不是最高代理分数，是因为决策目标为当前任务的预测效果。$\hat r_j$须来自同一损失、划分协议和终点，例如预算匹配的目标试训；LEEP与LogME可以安排试训顺序，却不能直接填成风险。若用代理预测终点风险，还需独立迁移记录建立并检验对应关系。已经花掉的共同搜索成本单独保留在累计账本中；若延迟也有硬要求，应在式~\eqref{eq:ket-reuse-single-choice}~中另加约束。

无迁移方案可行且预计风险严格低于所有可行来源时，选择结果为$j=0$，允许不用来源。如果整个可行集合为空，应报告没有可执行方案；无迁移不可行本身不意味着其他来源也不可行。若风险相同，按预先规定的成本等次级规则打破平局。优先节省标签时，可改为在效果约束下最小化标签预算。效果和成本同时重要时，只有存在另一方案效果不低、成本不高且至少一项严格改善，才剔除当前方案；其余候选再按应用要求取舍。

构造一个接口与预算反例：A的LEEP最高，却只提供远程概率查询，不能满足本地参数微调和离线部署。B预计准确率90\%、占400 MB，C预计89.8\%、占80 MB；部署上限100 MB时，B也被排除，C才能与其他可行候选及无迁移方案比较。能够计算评分不意味着能够执行最终复用，因而要在筛选前分别检查两种接口要求。

若放宽存储上限使B、C均可部署，二者0.2个百分点的预计差异还需与重复运行和样本估计的不确定性比较。证据不足以区分时，可以增加配对试训，检查关键病害类别，或暂时保留两者。硬预算要求立即决策时，则报告尚不能确认风险次序，并按预设的较低成本等规则选择。事后更换次级标准会把偏好混入结果，不增加新的效果证据。

选定一个来源后，还可以缩小使用范围：排除标签标准冲突的样本，取保留目标纹理的中间层，或只采用定位模块而不用旧分类头。记录中应写明模型版本、层名、样本纳入规则和模块接口，使证据对应到实际对象。缩小范围也有代价，删样本可能削弱覆盖，取早期层可能增加目标端重建高层结构的学习。第~\ref{chap:ket-transfer-adaptation}~章说明冻结、重训与新增任务参数的实施；这里先确定所用对象与选择依据。

\subsubsection{多来源知识的选择}
\label{sec:ket-reuse-multiple-sources}

一个来源覆盖不足时，我们可以检查是否有第二个来源补上缺口。但两个高准确率模型可能在同一批照片上出错，也可能因标签规范不同、概率含义不一致而发生冲突。因此，多来源选择的比较对象是来源集合按指定规则共同使用的结果。加入一个来源增加了多少目标收益，必须相对于已有集合计算，不能把单项成绩相加。

Taskonomy将来源任务的编码器分别训练并冻结，再以较少目标数据训练低容量读取器，测量有向迁移关系。多个来源表示共同进入读取器的高阶迁移，则用于检查组合互补。论文在各目标内部归一化这些关系，并用二元整数规划在监督预算下选择来源和迁移路径。这里的关系绑定任务集合、编码器、读取器及预算，不是任务之间天然固定的距离，也不描述共同训练时的梯度影响\scite{ket-zamir2018-taskonomy}

为说明多个目标怎样共用有限来源预算，先列出可实施路径$e=(J_e,k_e)$：来源集合$J_e$按指定组合方式服务目标$k_e$。用$\hat g_e$估计该路径相对目标基线的收益，不同目标先换算到约定效用尺度，$w_k\geq0$为目标权重。来源选中标志$u_j$和路径选中标志$v_e$均取0或1；$c_j,d_e\geq0$分别为来源成本与路径特有成本，按同一预算单位计量。我们要在总预算$B$内最大化加权收益，同时保证路径所需来源已经选中、每个目标有且仅有一条路径：
\begin{equation}
  \begin{aligned}
    \max_{\bm u,\bm v}\quad
      &\sum_e w_{k_e}\hat g_ev_e\\
    \text{满足}\quad
      &v_e\leq u_j \quad(j\in J_e),\\
      &\sum_{e:k_e=k}v_e=1 \quad(\text{每个目标 }k),\\
      &\sum_jc_ju_j+\sum_ed_ev_e\leq B.
  \end{aligned}
  \label{eq:ket-reuse-path-choice}
\end{equation}
第一组约束使选中的路径能取得全部来源，第二组为每个目标指定一条路径，第三组使共享来源只付一次成本、路径操作分别计费。每个目标可加入$J_e=\varnothing$的无来源路径，其收益为0，必要训练成本仍计入$d_e$，避免强迫迁移。式~\eqref{eq:ket-reuse-path-choice}~是教学成本扩展模型，不逐项复现Taskonomy约束。组合收益须来自组合试训或有依据的估计；即使准确求解整数规划，最优性也只针对列出的路径、估计和预算，不保证真实效果全局最优。

用表~\ref{tab:ket-reuse-complementarity}~考察按单项排名选择会遗漏什么。构造三个来源与两个等权目标：$\mathcal T_1$识别物种，$\mathcal T_2$判断病害，各有100个等权输入组、每组一个验证对象。组号表示输入可观察属性，而非样本行号；两项无迁移对照都在组1至20出错。这固定了计算收益的目标、分母和基准。

再规定组合计算规则。各来源模块在指定组上发出可用标志并修正基线，其他组不介入；标志只看组属性，不读真实标签。教学上安排触发的修正全部正确，重叠修正也一致。因此，任一模块提供修正时采用它，否则保留基线，组合修正范围恰为集合并集。这些条件使联合收益可以复算，并不是从各模块准确率推导出的真实集成性质。

具体修正集合如下：A对两目标均修正$\{1,\ldots,6\}$；B对$\mathcal T_1$修正$\{1,\ldots,8\}$，对$\mathcal T_2$修正$\{9\}$；C对$\mathcal T_1$修正$\{9\}$，对$\mathcal T_2$修正$\{1,\ldots,8\}$。每个来源加载与适应花1单位，组合额外成本设为零，总预算2单位。基线已完成，其共同成本不列在本表的新增成本中。表中每个组合按集合并集计数，来源A、B、C不代表真实模型。

\begin{table*}[htbp]
  \centering
  \small
  \begin{tabularx}{\linewidth}{lrrrrX}
    \toprule
    来源集合 & $\mathcal T_1$收益 & $\mathcal T_2$收益 & 平均收益 & 总成本 & 组合范围与预算状态 \\
    \midrule
    $\varnothing$ & 0 & 0 & 0 & 0 & 无迁移对照，各错20张 \\
    $\{A\}$ & 6 & 6 & 6 & 1 & 两个目标都修正组1至6 \\
    $\{B\}$ & 8 & 1 & 4.5 & 1 & 偏向物种目标 \\
    $\{C\}$ & 1 & 8 & 4.5 & 1 & 偏向病害目标 \\
    $\{A,B\}$ & 8 & 7 & 7.5 & 2 & 物种修正大部重叠 \\
    $\{A,C\}$ & 7 & 8 & 7.5 & 2 & 病害修正大部重叠 \\
    $\{B,C\}$ & 9 & 9 & 9 & 2 & 两个目标分别补上组9 \\
    $\{A,B,C\}$ & 9 & 9 & 9 & 3 & 超出预算，且A不增加覆盖 \\
    \bottomrule
  \end{tabularx}
  \caption{预算为2时的来源互补反例。收益是相对各目标20\%基线错误率的准确率增量，单位为百分点；成本只计新增来源加载与适应。全部组合由教学修正规则计算，实际模型须验证组合本身，不能由单项收益相加得到。}
  \label{tab:ket-reuse-complementarity}
\end{table*}

按单项平均收益排序先选A，再选B或C，联合平均收益只有7.5个百分点。枚举预算内集合则选$\{B,C\}$，两目标各提高9个百分点，错误数由20降为11。重复信息解释了这个差异：$\{A,B\}$在$\mathcal T_1$的修正并集只有$\{1,\ldots,8\}$，6与8不能加成14；在$\mathcal T_2$的并集为$\{1,\ldots,6,9\}$，计7。B、C单项均不如A，却分别补上另一个目标的组9。

若每步都计算实际边际收益，即加入新来源后相对当前集合增加的收益，第一步仍选A，第二步也只能达到7.5，未取得最优组合。并集覆盖有递减边际性质：已有修正范围越大，新来源还能补上的组越少；但这一性质本身也不保证逐步选择取得精确最优。真实模型加入来源还可能产生冲突、占用共享容量或改变路由，使原来正确的预测变错，连收益不下降都无法保证。只有明确集合函数性质和成本条件后才能讨论近似保证。小候选集可枚举，大候选集则先保留不同覆盖来源，再分配组合验证预算。

这个反例支持按组合而非单项排名选择，却没有给出真实模型的融合算法。第~\ref{chap:ket-integration}~章说明模块共同计算、预测融合、参数合并与蒸馏。若来源任务还要共同训练，来源表示本身会继续变化，应由第~\ref{chap:ket-sharing}~章分析共享，不再沿用这里固定来源、固定修正规则的组合成绩。

\subsection{选择结果的检验与调整}
\label{sec:ket-reuse-validation}

选中候选以后，继续适应得到的结果可能否定先前估计。应按预设口径核对目标效果、阈值到达量、实际支出和保留要求，并重点检查此前登记的缺口。整体准确率提高但新增物种全部被归入旧类，不能证明新增类别要求已满足；适应时间缩短但搜索支出更大，也不能证明全过程节省计算。检验所选方案要回到原先承诺的对象与指标。

风险高于对照时，还需区分原因。多个受控读取器都无法恢复目标标签，可以提示固定表示缺信息，但一般不能由有限探测证明所有读取方式无效。代理高分来源在匹配预算的试训中落后，说明该评分在此候选集发生误排；同一来源只改变学习率、冻结范围或标签映射后恢复效果，则提示适应设置起了作用。应围绕这些假设设计只改变相关因素的对照，不能从失败一次就归结为来源质量差。

据此可以重新选择来源，缩小到有支持的层或样本，隔离冲突能力，或停止迁移返回对照。选择不用来源也是受约束决策的允许结果。每次调整都要登记新增目标标签、无标签输入、试训轮次、版本和累计支出；只把同一批分数换一种排序，没有取得新的收益证据。

图~\ref{fig:ket-reuse-selection-loop}~把这些决策连成开发闭环。先由候选与访问条件确定可以取得的证据，估计可以直接采用试训，不要求必经代理评分。受约束选择分为单来源与多来源，再对所选方案继续试训、检查效果和成本。满足要求时锁定方案；不满足且仍有预算和可检验问题时调整；无可行候选、预算耗尽或继续试训不能解决关键不确定性时终止。停止规则应在开发前确定，不能无限寻找一个恰好满意的验证分数。

\begin{figure*}[htbp]
  \centering
  % Source file only. Development control flow is solid; evaluation data flow is dashed.
% All centerline connectors use shared x/y coordinates and explicit node anchors.
% Final evaluation has no return edge to development.
\begin{tikzpicture}[
  x=1cm,y=1cm,font=\figurefont\small,
  box/.style={
    rectangle,draw=BookInk,fill=BookPaper,line width=0.8pt,
    text width=42mm,minimum height=9mm,
    align=center,inner xsep=2mm,inner ysep=1.5mm
  },
  branch/.style={box,text width=28mm,fill=FigureModel!10},
  flow/.style={
    draw=BookInk,line width=0.9pt,rounded corners=1.2mm,
    -{Latex[length=2mm]}
  },
  edge label/.style={
    font=\figurefont\footnotesize,text=BookInk,inner sep=1mm
  }
]
  \node[box,fill=FigureInput!10,text width=50mm] (input)
    at (6,0) {候选集合与访问条件};
  \node[box] (score) at (6,-1.45) {候选价值估计\\固定评分或预算试训};
  \node[box] (choose) at (6,-2.9) {预算约束下选择};
  \node[branch] (single) at (3.2,-4.4) {单来源\\模型、层或子集};
  \node[branch] (multiple) at (8.8,-4.4) {多来源\\组合与路径};
  \node[box,fill=FigureModel!10] (trial)
    at (6,-5.95) {目标试训};
  \node[box] (review) at (6,-7.4) {判断：效果、成本\\与保留要求是否满足};
  \node[box,fill=FigureOutput!12] (lock)
    at (6,-9.3) {接受并锁定方案};
  \node[box,text width=33mm,fill=BookAmber!10] (stop)
    at (12.65,-5.95) {终止或回到对照};
  \node[box,text width=33mm,dashed,fill=white] (testdata)
    at (12.65,-7.7) {独立最终评价样本\\开发期间封存};
  \node[box,text width=33mm,fill=FigureOutput!12] (evaluate)
    at (12.65,-9.3) {最终评价与报告};

  \draw[flow] (input.south) -- (score.north);
  \draw[flow] (score.south) -- (choose.north);
  \draw[BookInk,line width=0.9pt] (choose.south) -- (6,-3.6);
  \draw[flow] (6,-3.6) -- (3.2,-3.6) -- (single.north);
  \draw[flow] (6,-3.6) -- (8.8,-3.6) -- (multiple.north);
  \draw[BookInk,line width=0.9pt,rounded corners=1.2mm]
    (single.south) -- (3.2,-5.25) -- (6,-5.25);
  \draw[BookInk,line width=0.9pt,rounded corners=1.2mm]
    (multiple.south) -- (8.8,-5.25) -- (6,-5.25);
  \draw[flow] (6,-5.25) -- (trial.north);
  \draw[flow] (trial.south) -- (review.north);
  \draw[flow] (review.south) -- (lock.north);
  \node[edge label,anchor=west] at (6.1,-8.55) {满足};
  \draw[flow] (lock.east) -- (evaluate.west);
  \draw[flow,dashed] (testdata.south) -- (evaluate.north);

  \draw[flow] (choose.east) -- (12.65,-2.9) -- (stop.north);
  \node[edge label,anchor=west] at (12.8,-4.35) {无可行候选};
  \draw[flow] (review.east) -- (9.55,-7.4)
    -- (9.55,-5.95) -- (stop.west);
  \node[edge label,anchor=west,align=left]
    at (9.65,-6.7) {不满足\\且不可继续};
  \draw[flow] (review.west) -- (0.4,-7.4)
    -- (0.4,-1.45) -- (score.west);
  \node[edge label,anchor=west,align=left]
    at (0.55,-6.45) {不满足且可继续\\调整候选或范围\\登记新增预算};
  \node[text=BookMuted,font=\figurefont\footnotesize]
    at (12.65,-10.15) {测试样本不返回选择环节};
\end{tikzpicture}

  \caption{知识选择的开发闭环。实线表示开发控制流程，单来源与多来源共享后续试训和要求检查；未满足且仍可继续时返回估计，并登记新增预算。虚线仅将封存的独立评价数据送入最终评价，没有返回开发的箭头。若据最终结果调整，原评价数据即已参与开发。}
  \label{fig:ket-reuse-selection-loop}
\end{figure*}

锁定时记录来源标识、复用范围、预处理、训练和停止规则，再执行独立最终评价。预算耗尽且证据仍不确定，应报告不足以确认收益；未检出负迁移不等于证明没有风险。最终评价发现新问题后可以开启下一轮研究，但若据此修改方案，原数据不再是下一轮独立证据，必须重新安排评价。第~\ref{chap:ket-evaluation}~章统一这些划分与统计协议。

\FloatBarrier
\section{本章小结}
\label{sec:ket-reuse-summary}

对于新地区的花卉目标，来源准确率和输入相似度分别描述旧能力与数据关系，都不能替代目标收益。我们先固定对照与预算，判断风险是否下降、达到指定效果是否少用标签、整个服务期限是否减少支出，再检查来源究竟提供了什么可用信息。同一个来源可以早期有利而终点有害，也可以适应便宜而全过程更贵；结论必须连同训练进度和成本边界保留。

选择证据的深度也有取舍。LEEP评价来源概率与目标标签的经验对应，LogME评价冻结特征在线性高斯代理中的标签证据；二者节省实际适应，却受接口、标签和模型假设限制。目标试训补充真实更新过程，代价是更多计算和选择数据使用。硬约束排除无法部署的高分来源，组合试验区分互补与重复，二者共同说明代理第一名并不必然是可执行方案中的最佳选择。

本章得到的选择是一份可以复核的条件化决定：采用哪个版本的哪些内容，依据什么输入、标签或任务关系，预计相对哪种对照改善哪个量，以及尚需支付什么成本。继续试训可以否定这一决定，预算与停止条件限制如何调整，独立最终评价检验锁定后的实际收益。来源缺信息、代理误排和适应设置不合适仍需不同对照辨认；证据不足或成本不合算时，无迁移方案应继续保留。
